\documentclass[11pt, a4paper]{article}

% --- Packages ---
\usepackage[utf8]{inputenc}
\usepackage{amsmath, amssymb, amsfonts, amsthm} % standard math packages
\usepackage[left=1in,right=1in,top=1in,bottom=1in]{geometry} % margins
\usepackage{enumitem} % for customizing lists (a), (b), etc.
\usepackage{fancyhdr} % for headers and footers
\usepackage{mathtools} % for better formatting of norms/abs values
\usepackage{tcolorbox} % styled problem boxes

% --- Header/Footer Setup ---
\pagestyle{fancy}
\fancyhf{}
\lhead{\textbf{MTL 107}}
\chead{Semester-II, 2025-2026}
\rhead{Tutorial Sheet 4}
\cfoot{\thepage}

% --- Custom Macros ---
\newcommand{\N}{\mathbb{N}}
\newcommand{\Z}{\mathbb{Z}}
\newcommand{\Q}{\mathbb{Q}}
\newcommand{\R}{\mathbb{R}}
\newcommand{\C}{\mathbb{C}}
\newcommand{\K}{\mathbb{K}} % For Field K
\newcommand{\eps}{\varepsilon}

% definition for Norm ||.|| and Absolute value |.|
\DeclarePairedDelimiter\abs{\lvert}{\rvert}
\DeclarePairedDelimiter\norm{\lVert}{\rVert}

% --- Environment Setup ---
\newenvironment{problem}{%
  \begin{tcolorbox}[
    colback=blue!4!white,
    colframe=blue!38!black,
    boxrule=0.95pt,
    arc=3pt,
    left=11pt,
    right=11pt,
    top=10pt,
    bottom=11pt,
    before skip=12pt,
    after skip=14pt,
  ]
  \noindent\textbf{\textcolor{blue!35!black}{Problem}}\par\nointerlineskip
  {\color{blue!42!black}\rule{\linewidth}{0.55pt}}\par\vspace{0.55em}%
}{%
  \end{tcolorbox}%
}
\newif\ifsolutionnewpage
\solutionnewpagetrue
\newenvironment{solution}
  {\begin{proof}[Solution]}
  {\end{proof}\ifsolutionnewpage\newpage\fi}

\title{\textbf{Tutorial Sheet 4 Solutions}}
\author{Prashant Tiwari}
\date{\today}

\begin{document}

\maketitle

% ----------------------------------------------------------------------
% INSTRUCTIONS:
% Copy the content provided by the AI for each problem and paste it below.
% ----------------------------------------------------------------------

% PASTE SOLUTION HERE


\begin{problem}
\begin{enumerate}[label=(\alph*)]
    \item Derive the three-point numerical differentiation formulas for approximating $f'(x)$ (midpoint and endpoint formulas) along with their corresponding error terms.
    \item Compare the resulting formulas in terms of their error terms and the number of function evaluations required for the approximation.
\end{enumerate}
\end{problem}

\begin{solution}
\textbf{Part (a): Derivation of Three-Point Formulas}

Let $x_0$ be the point where we want to evaluate the derivative, and let $h > 0$ be the step size. We use Taylor series expansions of $f(x)$ around $x_0$.

\vspace{0.2cm}
\textit{1. Three-Point Endpoint Formula (Forward Difference):} \\
Consider the nodes $x_0, x_0+h,$ and $x_0+2h$. We expand $f(x_0+h)$ and $f(x_0+2h)$ around $x_0$:
\begin{align}
f(x_0+h) &= f(x_0) + hf'(x_0) + \frac{h^2}{2}f''(x_0) + \frac{h^3}{6}f'''(\xi_1) \label{eq:e1} \\
f(x_0+2h) &= f(x_0) + 2hf'(x_0) + 2h^2f''(x_0) + \frac{8h^3}{6}f'''(\xi_2) \label{eq:e2}
\end{align}
where $\xi_1 \in (x_0, x_0+h)$ and $\xi_2 \in (x_0, x_0+2h)$.
To eliminate the $f''(x_0)$ term, we multiply equation \eqref{eq:e1} by 4 and subtract equation \eqref{eq:e2}:
\begin{equation*}
4f(x_0+h) - f(x_0+2h) = 3f(x_0) + 2hf'(x_0) + \frac{2h^3}{3}f'''(\xi_1) - \frac{4h^3}{3}f'''(\xi_2)
\end{equation*}
Solving for $f'(x_0)$, we get:
\begin{equation*}
f'(x_0) = \frac{-3f(x_0) + 4f(x_0+h) - f(x_0+2h)}{2h} - \frac{h^2}{3}\left[ 2f'''(\xi_2) - f'''(\xi_1) \right]
\end{equation*}
Using the Intermediate Value Theorem for continuous functions, the error terms can be combined into a single term evaluated at some $\xi \in (x_0, x_0+2h)$:
\begin{equation}
f'(x_0) = \frac{-3f(x_0) + 4f(x_0+h) - f(x_0+2h)}{2h} + \frac{h^2}{3}f'''(\xi)
\end{equation}
This is the three-point endpoint formula. The error term is $\frac{h^2}{3}f'''(\xi)$.

\vspace{0.2cm}
\textit{2. Three-Point Midpoint Formula (Central Difference):} \\
Consider the nodes $x_0-h, x_0,$ and $x_0+h$. We expand $f(x_0+h)$ and $f(x_0-h)$ around $x_0$:
\begin{align}
f(x_0+h) &= f(x_0) + hf'(x_0) + \frac{h^2}{2}f''(x_0) + \frac{h^3}{6}f'''(\xi_1) \label{eq:m1} \\
f(x_0-h) &= f(x_0) - hf'(x_0) + \frac{h^2}{2}f''(x_0) - \frac{h^3}{6}f'''(\xi_2) \label{eq:m2}
\end{align}
where $\xi_1 \in (x_0, x_0+h)$ and $\xi_2 \in (x_0-h, x_0)$.
Subtracting equation \eqref{eq:m2} from equation \eqref{eq:m1} perfectly eliminates both the $f(x_0)$ and $f''(x_0)$ terms:
\begin{equation*}
f(x_0+h) - f(x_0-h) = 2hf'(x_0) + \frac{h^3}{6}\left[f'''(\xi_1) + f'''(\xi_2)\right]
\end{equation*}
Solving for $f'(x_0)$ and combining the error terms via the Intermediate Value Theorem yields:
\begin{equation}
f'(x_0) = \frac{f(x_0+h) - f(x_0-h)}{2h} - \frac{h^2}{6}f'''(\xi)
\end{equation}
where $\xi \in (x_0-h, x_0+h)$. This is the three-point midpoint formula. The error term is $-\frac{h^2}{6}f'''(\xi)$.

\vspace{0.5cm}
\textbf{Part (b): Comparison}

\begin{itemize}
    \item \textbf{Error Term:} Both formulas are of order $\mathcal{O}(h^2)$. However, the coefficient of the error term in the midpoint formula ($\frac{1}{6}$) is half the magnitude of the coefficient in the endpoint formula ($\frac{1}{3}$). Consequently, the midpoint formula generally provides a significantly more accurate approximation.
    \item \textbf{Function Evaluations:} To compute the derivative at $x_0$, the midpoint formula requires evaluating the function at \textbf{2} points ($x_0+h$ and $x_0-h$). The endpoint formula requires evaluating the function at \textbf{3} points ($x_0$, $x_0+h$, and $x_0+2h$).
    \item \textbf{Conclusion:} The midpoint formula is superior because it yields a smaller theoretical truncation error and requires fewer function evaluations. The endpoint formula is typically reserved for boundaries (the first or last point in a given dataset) where centered points are not available.
\end{itemize}
\end{solution}
    


\begin{problem}
    The following data are taken from the function $f(x) = x \ln x$.
    
    \begin{center}
    \begin{tabular}{c | c | c}
    $x$ & $f(x)$ & $f'(x)$ \\
    \hline
    8.1 & 16.94410 & \\
    8.3 & 17.56492 & \\
    8.5 & 18.19056 & \\
    8.7 & 18.82091 &
    \end{tabular}
    \end{center}
    
    \begin{enumerate}[label=(\alph*)]
        \item Use the most accurate three-point numerical differentiation formulas to approximate $f'(x)$ at each listed value of $x$.
        \item Determine appropriate error bounds using the error formulas for the three-point methods, and then compute the actual errors in the approximations obtained in part (a).
    \end{enumerate}
    \end{problem}
    
    \begin{solution}
    The given function is $f(x) = x \ln x$. The uniform step size is $h = 0.2$. The most accurate 3-point formulas rely on the midpoint (central difference) formula for interior points, and the endpoint (forward/backward difference) formula for the boundaries.
    
    \vspace{0.3cm}
    \textbf{Part (a): Approximating $f'(x)$}
    
    \textbf{1. For $x = 8.1$ (Left Boundary):} \\
    We use the 3-point forward endpoint formula:
    \begin{equation*}
    f'(x_0) \approx \frac{-3f(x_0) + 4f(x_1) - f(x_2)}{2h}
    \end{equation*}
    \begin{align*}
    f'(8.1) &\approx \frac{-3(16.94410) + 4(17.56492) - 18.19056}{2(0.2)} \\
    &\approx \frac{-50.83230 + 70.25968 - 18.19056}{0.4} = \frac{1.23682}{0.4} = \mathbf{3.09205}
    \end{align*}
    
    \textbf{2. For $x = 8.3$ (Interior):} \\
    We use the 3-point midpoint formula:
    \begin{equation*}
    f'(x_1) \approx \frac{f(x_2) - f(x_0)}{2h}
    \end{equation*}
    \begin{align*}
    f'(8.3) &\approx \frac{18.19056 - 16.94410}{2(0.2)} = \frac{1.24646}{0.4} = \mathbf{3.11615}
    \end{align*}
    
    \textbf{3. For $x = 8.5$ (Interior):} \\
    We use the 3-point midpoint formula:
    \begin{align*}
    f'(8.5) &\approx \frac{f(8.7) - f(8.3)}{2(0.2)} = \frac{18.82091 - 17.56492}{0.4} = \frac{1.25599}{0.4} = \mathbf{3.139975}
    \end{align*}
    
    \textbf{4. For $x = 8.7$ (Right Boundary):} \\
    We use the 3-point backward endpoint formula:
    \begin{equation*}
    f'(x_n) \approx \frac{3f(x_n) - 4f(x_{n-1}) + f(x_{n-2})}{2h}
    \end{equation*}
    \begin{align*}
    f'(8.7) &\approx \frac{3(18.82091) - 4(18.19056) + 17.56492}{2(0.2)} \\
    &\approx \frac{56.46273 - 72.76224 + 17.56492}{0.4} = \frac{1.26541}{0.4} = \mathbf{3.163525}
    \end{align*}
    
    \vspace{0.3cm}
    \textbf{Part (b): Error Bounds and Actual Errors}
    
    To find the theoretical bounds, we need the third derivative of $f(x) = x \ln x$.
    \begin{align*}
    f'(x) &= \ln x + 1 \\
    f''(x) &= \frac{1}{x} \\
    f'''(x) &= -\frac{1}{x^2}
    \end{align*}
    Since $\abs{f'''(x)} = \frac{1}{x^2}$ is a strictly decreasing function for $x > 0$, its maximum absolute value on any interval $[a, b]$ will occur at the left endpoint, $x = a$.
    
    \vspace{0.2cm}
    \textbf{Error Bounds:}
    \begin{itemize}
        \item \textbf{For $x = 8.1$ (Endpoint):} The interval for $\xi$ is $[8.1, 8.5]$.
        $$ \text{Bound} = \frac{h^2}{3} \max_{\xi \in [8.1, 8.5]} \abs[\Big]{-\frac{1}{\xi^2}} = \frac{0.04}{3} \left( \frac{1}{8.1^2} \right) \approx \mathbf{2.032 \times 10^{-4}} $$
    
        \item \textbf{For $x = 8.3$ (Midpoint):} The interval for $\xi$ is $[8.1, 8.5]$.
        $$ \text{Bound} = \frac{h^2}{6} \max_{\xi \in [8.1, 8.5]} \abs[\Big]{-\frac{1}{\xi^2}} = \frac{0.04}{6} \left( \frac{1}{8.1^2} \right) \approx \mathbf{1.016 \times 10^{-4}} $$
    
        \item \textbf{For $x = 8.5$ (Midpoint):} The interval for $\xi$ is $[8.3, 8.7]$.
        $$ \text{Bound} = \frac{h^2}{6} \max_{\xi \in [8.3, 8.7]} \abs[\Big]{-\frac{1}{\xi^2}} = \frac{0.04}{6} \left( \frac{1}{8.3^2} \right) \approx \mathbf{9.677 \times 10^{-5}} $$
    
        \item \textbf{For $x = 8.7$ (Endpoint):} The interval for $\xi$ is $[8.3, 8.7]$.
        $$ \text{Bound} = \frac{h^2}{3} \max_{\xi \in [8.3, 8.7]} \abs[\Big]{-\frac{1}{\xi^2}} = \frac{0.04}{3} \left( \frac{1}{8.3^2} \right) \approx \mathbf{1.935 \times 10^{-4}} $$
    \end{itemize}
    
    \vspace{0.2cm}
    \textbf{Actual Errors:} \\
    The true values are computed using $f'(x) = \ln(x) + 1$:
    \begin{itemize}
        \item $f'(8.1) = \ln(8.1) + 1 \approx 3.091864$
        \item $f'(8.3) = \ln(8.3) + 1 \approx 3.116256$
        \item $f'(8.5) = \ln(8.5) + 1 \approx 3.140066$
        \item $f'(8.7) = \ln(8.7) + 1 \approx 3.163323$
    \end{itemize}
    
    We calculate the absolute actual error as $E = \abs{\text{True} - \text{Approximation}}$:
    \begin{itemize}
        \item \textbf{At $x = 8.1$:} $\abs{3.091864 - 3.092050} = \mathbf{1.86 \times 10^{-4}}$ (Complies with bound $\le 2.032 \times 10^{-4}$)
        \item \textbf{At $x = 8.3$:} $\abs{3.116256 - 3.116150} = \mathbf{1.06 \times 10^{-4}}$ (\textit{Slightly exceeds bound $\le 1.016 \times 10^{-4}$})
        \item \textbf{At $x = 8.5$:} $\abs{3.140066 - 3.139975} = \mathbf{0.91 \times 10^{-4}}$ (Complies with bound $\le 0.967 \times 10^{-4}$)
        \item \textbf{At $x = 8.7$:} $\abs{3.163323 - 3.163525} = \mathbf{2.02 \times 10^{-4}}$ (\textit{Slightly exceeds bound $\le 1.935 \times 10^{-4}$})
    \end{itemize}
    
    \textit{Note: The slight discrepancies where the actual error exceeds the theoretical bound for $x=8.3$ and $x=8.7$ are due to the accumulated round-off errors present in the provided five-decimal-place data table, as the theoretical bounds only account for pure truncation error.}
    \end{solution}

    

\begin{problem}
    The following data are taken from the function $f(x) = e^{x/3} + x^2$.
    
    \begin{center}
    \begin{tabular}{c | c | c}
    $x$ & $f(x)$ & $f'(x)$ \\
    \hline
    -3.0 & 9.367879 & \\
    -2.8 & 8.233241 & \\
    -2.6 & 7.180350 & \\
    -2.4 & 6.209329 & \\
    -2.2 & 5.320305 & \\
    -2.0 & 4.513417 &
    \end{tabular}
    \end{center}
    
    \begin{enumerate}[label=(\alph*)]
        \item Use the most accurate five-point numerical differentiation formulas to approximate $f'(x)$ at each listed value of $x$.
        \item Determine appropriate error bounds using the error formulas for the five-point methods, and then compute the actual errors in the approximations obtained in part (a).
    \end{enumerate}
    \end{problem}
    
    \begin{solution}
    The given function is $f(x) = e^{x/3} + x^2$. The step size is $h = 0.2$. We are given six data points $x_0, x_1, x_2, x_3, x_4, x_5$.
    
    \vspace{0.3cm}
    \textbf{Part (a): Approximating $f'(x)$}
    
    We use the 5-point midpoint formula where centered data is available, and the 5-point endpoint formulas for the boundaries.
    \begin{itemize}
        \item \textbf{Midpoint Formula (for $x_2, x_3$):}
        $$ f'(x) \approx \frac{1}{12h} \left[ f(x-2h) - 8f(x-h) + 8f(x+h) - f(x+2h) \right] $$
        \item \textbf{Endpoint Forward Formula (for $x_0, x_1$):}
        $$ f'(x) \approx \frac{1}{12h} \left[ -25f(x) + 48f(x+h) - 36f(x+2h) + 16f(x+3h) - 3f(x+4h) \right] $$
        \item \textbf{Endpoint Backward Formula (for $x_4, x_5$):}
        $$ f'(x) \approx \frac{1}{12h} \left[ 3f(x-4h) - 16f(x-3h) + 36f(x-2h) - 48f(x-h) + 25f(x) \right] $$
    \end{itemize}
    
    Applying these formulas ($12h = 2.4$):
    \begin{itemize}
        \item \textbf{For $x = -3.0$ (Forward using $x_0$ to $x_4$):}
        \begin{align*}
        f'(-3.0) &\approx \frac{1}{2.4} \left[ -25(9.367879) + 48(8.233241) - 36(7.180350) + 16(6.209329) - 3(5.320305) \right] \\
        &\approx \mathbf{-5.877354}
        \end{align*}
    
        \item \textbf{For $x = -2.8$ (Forward using $x_1$ to $x_5$):}
        \begin{align*}
        f'(-2.8) &\approx \frac{1}{2.4} \left[ -25(8.233241) + 48(7.180350) - 36(6.209329) + 16(5.320305) - 3(4.513417) \right] \\
        &\approx \mathbf{-5.468962}
        \end{align*}
    
        \item \textbf{For $x = -2.6$ (Midpoint using $x_0$ to $x_4$):}
        \begin{align*}
        f'(-2.6) &\approx \frac{1}{2.4} \left[ 9.367879 - 8(8.233241) + 8(6.209329) - 5.320305 \right] \approx \mathbf{-5.059871}
        \end{align*}
    
        \item \textbf{For $x = -2.4$ (Midpoint using $x_1$ to $x_5$):}
        \begin{align*}
        f'(-2.4) &\approx \frac{1}{2.4} \left[ 8.233241 - 8(7.180350) + 8(5.320305) - 4.513417 \right] \approx \mathbf{-4.650381}
        \end{align*}
    
        \item \textbf{For $x = -2.2$ (Backward using $x_0$ to $x_4$):}
        \begin{align*}
        f'(-2.2) &\approx \frac{1}{2.4} \left[ 3(9.367879) - 16(8.233241) + 36(7.180350) - 48(6.209329) + 25(5.320305) \right] \\
        &\approx \mathbf{-4.240562}
        \end{align*}
    
        \item \textbf{For $x = -2.0$ (Backward using $x_1$ to $x_5$):}
        \begin{align*}
        f'(-2.0) &\approx \frac{1}{2.4} \left[ 3(8.233241) - 16(7.180350) + 36(6.209329) - 48(5.320305) + 25(4.513417) \right] \\
        &\approx \mathbf{-3.828859}
        \end{align*}
    \end{itemize}
    
    \vspace{0.3cm}
    \textbf{Part (b): Error Bounds and Actual Errors}
    
    First, we need the fifth derivative of $f(x)$:
    \begin{align*}
    f'(x) &= \frac{1}{3}e^{x/3} + 2x \\
    f''(x) &= \frac{1}{9}e^{x/3} + 2 \\
    f'''(x) &= \frac{1}{27}e^{x/3} \\
    f^{(4)}(x) &= \frac{1}{81}e^{x/3} \\
    f^{(5)}(x) &= \frac{1}{243}e^{x/3}
    \end{align*}
    Since $e^{x/3}$ is an increasing function, its maximum absolute value on the interval $[-3.0, -2.0]$ occurs at the rightmost endpoint, $x = -2.0$.
    $$ \max_{\xi \in [-3.0, -2.0]} \abs[\Big]{f^{(5)}(\xi)} = \frac{1}{243}e^{-2.0/3} \approx \frac{1}{243}(0.513417) \approx 0.0021128 $$
    
    \textbf{Theoretical Error Bounds:}
    \begin{itemize}
        \item \textbf{Endpoint Formula Bound} (for $x = -3.0, -2.8, -2.2, -2.0$):
        $$ \frac{h^4}{5} \abs[\Big]{f^{(5)}(\xi)} \le \frac{(0.2)^4}{5}(0.0021128) = \frac{0.0016}{5}(0.0021128) \approx \mathbf{6.76 \times 10^{-7}} $$
        \item \textbf{Midpoint Formula Bound} (for $x = -2.6, -2.4$):
        $$ \frac{h^4}{30} \abs[\Big]{f^{(5)}(\xi)} \le \frac{(0.2)^4}{30}(0.0021128) = \frac{0.0016}{30}(0.0021128) \approx \mathbf{1.13 \times 10^{-7}} $$
    \end{itemize}
    
    \textbf{Actual Errors:}
    We evaluate the exact derivative $f'(x) = \frac{1}{3}e^{x/3} + 2x$:
    \begin{itemize}
        \item $f'(-3.0) = \frac{1}{3}e^{-1.0} - 6 \approx -5.877373$
        \item $f'(-2.8) = \frac{1}{3}e^{-2.8/3} - 5.6 \approx -5.468965$
        \item $f'(-2.6) = \frac{1}{3}e^{-2.6/3} - 5.2 \approx -5.059871$
        \item $f'(-2.4) = \frac{1}{3}e^{-0.8} - 4.8 \approx -4.650223$
        \item $f'(-2.2) = \frac{1}{3}e^{-2.2/3} - 4.4 \approx -4.239850$
        \item $f'(-2.0) = \frac{1}{3}e^{-2.0/3} - 4.0 \approx -3.828861$
    \end{itemize}
    
    Computing the absolute actual errors $\abs{\text{True} - \text{Approximation}}$:
    \begin{itemize}
        \item \textbf{At $x = -3.0$:} $\abs{-5.877373 - (-5.877354)} = \mathbf{1.9 \times 10^{-5}}$
        \item \textbf{At $x = -2.8$:} $\abs{-5.468965 - (-5.468962)} = \mathbf{3.0 \times 10^{-6}}$
        \item \textbf{At $x = -2.6$:} $\abs{-5.059871 - (-5.059871)} = \mathbf{0.0 \times 10^{-6}}$
        \item \textbf{At $x = -2.4$:} $\abs{-4.650223 - (-4.650381)} = \mathbf{1.58 \times 10^{-4}}$
        \item \textbf{At $x = -2.2$:} $\abs{-4.239850 - (-4.240562)} = \mathbf{7.12 \times 10^{-4}}$
        \item \textbf{At $x = -2.0$:} $\abs{-3.828861 - (-3.828859)} = \mathbf{2.0 \times 10^{-6}}$
    \end{itemize}
    \textit{Note: As before, actual errors can exceed truncation theoretical bounds due to the round-off error present in the provided 6-decimal point starting data.}
    \end{solution}

    \begin{problem}
        Consider the three-point midpoint formula for approximating the derivative
        \begin{equation*}
        f'(x_0) = \frac{f(x_0+h) - f(x_0-h)}{2h} - \frac{h^2}{6}f^{(3)}(\xi_1).
        \end{equation*}
        Suppose that in evaluating $f(x_0+h)$ and $f(x_0-h)$ we encounter round-off errors $e(x_0+h)$ and $e(x_0-h)$. The computed values $\overline{f}(x_0+h)$ and $\overline{f}(x_0-h)$ are related to the true values by
        \begin{align*}
        f(x_0+h) &= \overline{f}(x_0+h) + e(x_0+h), \\
        f(x_0-h) &= \overline{f}(x_0-h) + e(x_0-h).
        \end{align*}
        \begin{enumerate}[label=(\alph*)]
            \item Derive an error bound for the approximation of $f'(x_0)$ obtained from the three-point midpoint formula when both truncation error and round-off errors are present. Assume $\abs{e(x_0+h)} \le \eps$, $\abs{e(x_0-h)} \le \eps$, and $\abs{f^{(3)}(x)} \le M$ for $x$ in the relevant interval.
            \item Based on the derived bound, discuss whether it is practical to take $h$ as small as possible. Determine an appropriate choice of $h$ that balances truncation and round-off errors.
        \end{enumerate}
        \end{problem}
        
        \begin{solution}
        \textbf{Part (a): Deriving the Error Bound}
        
        Let the approximation of the derivative using the computed (numerical) values be denoted by $\tilde{f}'(x_0)$:
        \begin{equation*}
        \tilde{f}'(x_0) = \frac{\overline{f}(x_0+h) - \overline{f}(x_0-h)}{2h}
        \end{equation*}
        The exact mathematical relationship, taking into account the truncation error, is:
        \begin{equation*}
        f'(x_0) = \frac{f(x_0+h) - f(x_0-h)}{2h} - \frac{h^2}{6}f^{(3)}(\xi_1)
        \end{equation*}
        Substituting the expressions for the true values $f(x_0+h)$ and $f(x_0-h)$ in terms of their computed values and round-off errors:
        \begin{align*}
        f'(x_0) &= \frac{[\overline{f}(x_0+h) + e(x_0+h)] - [\overline{f}(x_0-h) + e(x_0-h)]}{2h} - \frac{h^2}{6}f^{(3)}(\xi_1) \\
        &= \frac{\overline{f}(x_0+h) - \overline{f}(x_0-h)}{2h} + \frac{e(x_0+h) - e(x_0-h)}{2h} - \frac{h^2}{6}f^{(3)}(\xi_1) \\
        &= \tilde{f}'(x_0) + \frac{e(x_0+h) - e(x_0-h)}{2h} - \frac{h^2}{6}f^{(3)}(\xi_1)
        \end{align*}
        The total error is the difference between the true derivative and the computed approximation:
        \begin{equation*}
        E(h) = f'(x_0) - \tilde{f}'(x_0) = \frac{e(x_0+h) - e(x_0-h)}{2h} - \frac{h^2}{6}f^{(3)}(\xi_1)
        \end{equation*}
        Taking the absolute value and applying the triangle inequality:
        \begin{equation*}
        \abs{E(h)} \le \frac{\abs{e(x_0+h)} + \abs{e(x_0-h)}}{2h} + \frac{h^2}{6}\abs{f^{(3)}(\xi_1)}
        \end{equation*}
        Given that $\abs{e(x_0+h)} \le \eps$, $\abs{e(x_0-h)} \le \eps$, and $\abs{f^{(3)}(\xi_1)} \le M$, we obtain the total error bound:
        \begin{equation}
        \abs{E(h)} \le \frac{\eps + \eps}{2h} + \frac{h^2}{6}M = \frac{\eps}{h} + \frac{Mh^2}{6}
        \end{equation}
        
        \vspace{0.3cm}
        \textbf{Part (b): Optimal Choice of $h$}
        
        Let the upper bound of the total error be denoted by $B(h)$:
        \begin{equation*}
        B(h) = \frac{\eps}{h} + \frac{M}{6}h^2
        \end{equation*}
        
        \textbf{Is it practical to take $h$ as small as possible?} \\
        No, it is not practical. As $h \to 0$, the truncation error portion $\left(\frac{M}{6}h^2\right)$ approaches zero, but the round-off error portion $\left(\frac{\eps}{h}\right)$ approaches infinity. Thus, making $h$ arbitrarily small will cause the total error to blow up because dividing the inherent machine precision limits ($\eps$) by a tiny number amplifies the error.
        
        \textbf{Determining the appropriate choice of $h$:} \\
        To find the step size $h$ that minimizes the maximum error bound $B(h)$, we take the derivative of $B(h)$ with respect to $h$ and set it to zero:
        \begin{equation*}
        B'(h) = -\frac{\eps}{h^2} + \frac{M}{3}h = 0
        \end{equation*}
        Solving for $h$:
        \begin{align*}
        \frac{M}{3}h &= \frac{\eps}{h^2} \\
        h^3 &= \frac{3\eps}{M} \\
        h &= \left( \frac{3\eps}{M} \right)^{1/3}
        \end{align*}
        To verify this is a minimum, we check the second derivative:
        \begin{equation*}
        B''(h) = \frac{2\eps}{h^3} + \frac{M}{3}
        \end{equation*}
        Since $\eps, M, h > 0$, it follows that $B''(h) > 0$, confirming that $h = \left( \frac{3\eps}{M} \right)^{1/3}$ yields the minimum error bound. This is the optimal step size $h$ that perfectly balances the truncation error against the machine round-off error.
        \end{solution}

        
\begin{problem}
    Use Heun's method of order two to approximate the solution of the initial-value problem
    \begin{equation*}
    y' = te^{3t} - 2y, \quad 0 \le t \le 1, \quad y(0) = 0
    \end{equation*}
    with step size $h = 0.5$.
    \end{problem}
    
    \begin{solution}
    The initial-value problem is defined by $f(t, y) = te^{3t} - 2y$, with initial conditions $t_0 = 0$, $w_0 = y(0) = 0$, and a step size of $h = 0.5$.
    We need to find approximations $w_1$ at $t_1 = 0.5$ and $w_2$ at $t_2 = 1.0$.
    
    Heun's method of order two (also known as the explicit trapezoidal rule) consists of a predictor step (Euler's method) and a corrector step:
    \begin{align*}
    \text{Predictor: } \tilde{w}_{i+1} &= w_i + h f(t_i, w_i) \\
    \text{Corrector: } w_{i+1} &= w_i + \frac{h}{2} \left[ f(t_i, w_i) + f(t_{i+1}, \tilde{w}_{i+1}) \right]
    \end{align*}
    
    \vspace{0.3cm}
    \textbf{Step 1: Approximation at $t_1 = 0.5$}
    \begin{itemize}
        \item Evaluate $f(t_0, w_0)$:
        $$ f(0, 0) = 0 \cdot e^{3(0)} - 2(0) = 0 $$
    
        \item Compute the predictor $\tilde{w}_1$:
        $$ \tilde{w}_1 = w_0 + h f(0, 0) = 0 + 0.5(0) = 0 $$
    
        \item Evaluate $f(t_1, \tilde{w}_1)$ at $t_1 = 0.5$:
        $$ f(0.5, 0) = 0.5 e^{3(0.5)} - 2(0) = 0.5 e^{1.5} \approx 2.240845 $$
    
        \item Compute the corrector $w_1$:
        $$ w_1 = w_0 + \frac{0.5}{2} [ f(0, 0) + f(0.5, 0) ] $$
        $$ w_1 \approx 0 + 0.25 [ 0 + 2.240845 ] = \mathbf{0.560211} $$
    \end{itemize}
    Thus, the approximation at $t = 0.5$ is $y(0.5) \approx 0.560211$.
    
    \vspace{0.3cm}
    \textbf{Step 2: Approximation at $t_2 = 1.0$}
    \begin{itemize}
        \item Evaluate $f(t_1, w_1)$:
        $$ f(0.5, 0.560211) = 0.5 e^{1.5} - 2(0.560211) \approx 2.240845 - 1.120422 = 1.120423 $$
    
        \item Compute the predictor $\tilde{w}_2$:
        $$ \tilde{w}_2 = w_1 + h f(t_1, w_1) \approx 0.560211 + 0.5(1.120423) = 1.120423 $$
    
        \item Evaluate $f(t_2, \tilde{w}_2)$ at $t_2 = 1.0$:
        $$ f(1.0, 1.120423) = 1.0 e^{3(1.0)} - 2(1.120423) \approx 20.085537 - 2.240846 = 17.844691 $$
    
        \item Compute the corrector $w_2$:
        $$ w_2 = w_1 + \frac{0.5}{2} [ f(t_1, w_1) + f(t_2, \tilde{w}_2) ] $$
        $$ w_2 \approx 0.560211 + 0.25 [ 1.120423 + 17.844691 ] $$
        $$ w_2 \approx 0.560211 + 0.25 [ 18.965114 ] \approx 0.560211 + 4.741279 = \mathbf{5.301490} $$
    \end{itemize}
    
    \textbf{Final Answer:} \\
    The approximate solutions are $w_1 \approx \mathbf{0.560211}$ at $t=0.5$, and $w_2 \approx \mathbf{5.301490}$ at $t=1.0$.
    \end{solution}

    
\begin{problem}
    Suppose $f$ is continuous and satisfies a Lipschitz condition with constant $L$ on
    $D = \{ (t,y) \mid a \le t \le b, -\infty < y < \infty \},$
    and suppose there exists a constant $M$ such that $|y''(t)| \le M$ for all $t \in [a,b]$, where $y(t)$ is the unique solution of the initial-value problem
    \begin{equation*}
    y' = f(t,y), \quad a \le t \le b, \quad y(a) = \alpha
    \end{equation*}
    \begin{enumerate}[label=(\alph*)]
        \item Let $w_0, w_1, ..., w_N$ be the approximations generated by Euler's method with step size $h$. Prove that for each $i=0,1,2,...,N$,
        \begin{equation*}
        |y(t_i) - w_i| \le \frac{hM}{2L}(e^{L(t_i-a)} - 1).
        \end{equation*}
        \item Suppose the approximations $u_0, u_1, ..., u_N$ are generated by the perturbed Euler method
        \begin{equation*}
        u_0 = \alpha + \delta_0, \quad u_{i+1} = u_i + hf(t_i, u_i) + \delta_{i+1}, \quad i=0,1,...,N-1
        \end{equation*}
        where $\delta_i$ represents the round-off error at each step and $|\delta_i| \le \delta$. Show that for each $i=0,1,...,N$;
        \begin{equation*}
        |y(t_i) - u_i| \le \frac{1}{L}\left(\frac{hM}{2} + \frac{\delta}{h}\right)[e^{L(t_i-a)} - 1] + |\delta_0|e^{L(t_i-a)}.
        \end{equation*}
        Using this bound, discuss the effect of round-off errors and determine an appropriate choice of the step size $h$.
    \end{enumerate}
    \end{problem}
    
    \begin{solution}
    To prove both parts, we will use the following well-known lemma for difference inequalities: \\
    \textbf{Lemma:} If $s > 0$ and $t > 0$, and a sequence $\{a_i\}_{i=0}^N$ satisfies $a_{i+1} \le (1+s)a_i + t$, then
    \begin{equation} \label{eq:lemma}
    a_i \le e^{is}a_0 + \frac{t}{s}(e^{is}-1)
    \end{equation}
    Also note that $is = i(hL) = (ih)L = L(t_i - a)$, since $t_i = a + ih$.
    
    \vspace{0.3cm}
    \textbf{Part (a): Error Bound for Standard Euler's Method}
    
    Using Taylor's Theorem, the exact solution $y(t)$ satisfies:
    \begin{equation} \label{eq:exact}
    y(t_{i+1}) = y(t_i) + h f(t_i, y(t_i)) + \frac{h^2}{2} y''(\xi_i)
    \end{equation}
    for some $\xi_i \in (t_i, t_{i+1})$. Euler's method generates approximations via:
    \begin{equation} \label{eq:euler}
    w_{i+1} = w_i + h f(t_i, w_i)
    \end{equation}
    Subtracting \eqref{eq:euler} from \eqref{eq:exact} gives the error equation:
    \begin{equation*}
    y(t_{i+1}) - w_{i+1} = y(t_i) - w_i + h [f(t_i, y(t_i)) - f(t_i, w_i)] + \frac{h^2}{2} y''(\xi_i)
    \end{equation*}
    Taking the absolute value and applying the triangle inequality:
    \begin{equation*}
    |y(t_{i+1}) - w_{i+1}| \le |y(t_i) - w_i| + h |f(t_i, y(t_i)) - f(t_i, w_i)| + \frac{h^2}{2} |y''(\xi_i)|
    \end{equation*}
    Since $f$ satisfies a Lipschitz condition with constant $L$, $|f(t_i, y(t_i)) - f(t_i, w_i)| \le L|y(t_i) - w_i|$. We are also given $|y''(t)| \le M$. Let $e_i = |y(t_i) - w_i|$. Substituting these bounds yields:
    \begin{equation*}
    e_{i+1} \le e_i + hL e_i + \frac{h^2 M}{2} = (1 + hL)e_i + \frac{h^2 M}{2}
    \end{equation*}
    Now we apply the Lemma \eqref{eq:lemma} with $a_i = e_i$, $s = hL$, and $t = \frac{h^2 M}{2}$.
    Since the initial condition is exact in the standard Euler method, $w_0 = \alpha = y(t_0)$, meaning $e_0 = 0$.
    \begin{align*}
    e_i &\le e^{i(hL)}(0) + \frac{h^2 M / 2}{hL} (e^{i(hL)} - 1) \\
    |y(t_i) - w_i| &\le \frac{hM}{2L} (e^{L(t_i - a)} - 1)
    \end{align*}
    This completes the proof for part (a).
    
    \vspace{0.3cm}
    \textbf{Part (b): Error Bound for Perturbed Euler's Method}
    
    The perturbed Euler method is given by:
    \begin{equation} \label{eq:pert}
    u_{i+1} = u_i + h f(t_i, u_i) + \delta_{i+1}
    \end{equation}
    Subtracting \eqref{eq:pert} from the exact Taylor expansion \eqref{eq:exact}:
    \begin{equation*}
    y(t_{i+1}) - u_{i+1} = y(t_i) - u_i + h [f(t_i, y(t_i)) - f(t_i, u_i)] + \frac{h^2}{2} y''(\xi_i) - \delta_{i+1}
    \end{equation*}
    Taking the absolute value, applying the triangle inequality, and using the Lipschitz condition and bounds $|y''(t)| \le M$ and $|\delta_i| \le \delta$:
    \begin{align*}
    |y(t_{i+1}) - u_{i+1}| &\le |y(t_i) - u_i| + h L |y(t_i) - u_i| + \frac{h^2}{2} M + \delta \\
    |y(t_{i+1}) - u_{i+1}| &\le (1 + hL)|y(t_i) - u_i| + \left( \frac{h^2 M}{2} + \delta \right)
    \end{align*}
    Let $E_i = |y(t_i) - u_i|$. Applying the Lemma \eqref{eq:lemma} with $a_i = E_i$, $s = hL$, and $t = \frac{h^2 M}{2} + \delta$:
    \begin{equation*}
    E_i \le e^{i(hL)} E_0 + \frac{\frac{h^2 M}{2} + \delta}{hL} (e^{i(hL)} - 1)
    \end{equation*}
    Here, the initial error is $E_0 = |y(t_0) - u_0| = |\alpha - (\alpha + \delta_0)| = |\delta_0|$. Replacing $i(hL)$ with $L(t_i - a)$:
    \begin{equation*}
    |y(t_i) - u_i| \le e^{L(t_i - a)}|\delta_0| + \frac{1}{L}\left( \frac{hM}{2} + \frac{\delta}{h} \right) [e^{L(t_i - a)} - 1]
    \end{equation*}
    This completes the proof for the perturbed error bound.
    
    \vspace{0.3cm}
    \textbf{Discussion on Round-off Errors and Step Size:}
    
    Let the term dependent on $h$ in the error bound be denoted by $E(h)$:
    \begin{equation*}
    E(h) = \frac{hM}{2} + \frac{\delta}{h}
    \end{equation*}
    \begin{itemize}
        \item \textbf{Effect of round-off errors:} As the step size $h$ decreases ($h \to 0$), the mathematical truncation error $\frac{hM}{2}$ approaches zero. However, the accumulated round-off error term $\frac{\delta}{h}$ grows without bound (approaching infinity). Therefore, unlike the idealized mathematical method, we cannot make the actual error arbitrarily small simply by shrinking $h$.
        \item \textbf{Appropriate choice of step size $h$:} To minimize the maximum error bound $E(h)$, we find its critical point by setting its derivative with respect to $h$ to zero:
        \begin{align*}
        E'(h) = \frac{M}{2} - \frac{\delta}{h^2} &= 0 \\
        \frac{\delta}{h^2} &= \frac{M}{2} \\
        h^2 &= \frac{2\delta}{M} \implies h = \sqrt{\frac{2\delta}{M}}
        \end{align*}
        Since $E''(h) = \frac{2\delta}{h^3} > 0$ for positive $h$, this critical point represents a minimum. Thus, the optimal step size that balances truncation error and round-off error is $h = \sqrt{\frac{2\delta}{M}}$.
    \end{itemize}
    \end{solution}

    
    \begin{problem}
        Use the Modified Euler method and the Runge-Kutta method of order four to approximate the solution of the initial-value problem
        \begin{equation*}
        y' = te^{3t} - 2y, \quad 0 \le t \le 1, \quad y(0) = 0
        \end{equation*}
        with step size $h = 0.5$. Compare the numerical results with the exact solution
        \begin{equation*}
        y(t) = \frac{1}{5}te^{3t} - \frac{1}{25}e^{3t} + \frac{1}{25}e^{-2t}.
        \end{equation*}
        \end{problem}
        
        \begin{solution}
        The initial-value problem is defined by $f(t, y) = te^{3t} - 2y$, with $t_0 = 0, w_0 = 0$, and step size $h = 0.5$. We need approximations at $t_1 = 0.5$ and $t_2 = 1.0$.
        
        \vspace{0.3cm}
        \textbf{1. Modified Euler Method} 
        
        The Modified Euler method (also known as the Midpoint Method) is given by:
        \begin{equation*}
        w_{i+1} = w_i + h f\left(t_i + \frac{h}{2}, w_i + \frac{h}{2}f(t_i, w_i)\right)
        \end{equation*}
        
        \textit{Step 1: Approximation at $t_1 = 0.5$}
        \begin{align*}
        f(t_0, w_0) &= f(0, 0) = 0 \\
        w_1 &= 0 + 0.5 f(0 + 0.25, 0 + 0.25(0)) \\
        &= 0.5 f(0.25, 0) = 0.5 \left[ 0.25e^{3(0.25)} - 2(0) \right] \\
        &= 0.5 [0.25e^{0.75}] \approx 0.5 [0.529250] = \mathbf{0.264625}
        \end{align*}
        
        \textit{Step 2: Approximation at $t_2 = 1.0$}
        \begin{align*}
        f(t_1, w_1) &= f(0.5, 0.264625) = 0.5e^{1.5} - 2(0.264625) \approx 2.240845 - 0.529250 = 1.711595 \\
        w_2 &= 0.264625 + 0.5 f(0.5 + 0.25, 0.264625 + 0.25(1.711595)) \\
        &= 0.264625 + 0.5 f(0.75, 0.692524) \\
        &= 0.264625 + 0.5 \left[ 0.75e^{3(0.75)} - 2(0.692524) \right] \\
        &\approx 0.264625 + 0.5 \left[ 7.115801 - 1.385048 \right] \approx 0.264625 + 0.5(5.730753) = \mathbf{3.130001}
        \end{align*}
        
        \vspace{0.3cm}
        \textbf{2. Runge-Kutta Method of Order Four (RK4)}
        
        The RK4 method is given by $w_{i+1} = w_i + \frac{1}{6}(k_1 + 2k_2 + 2k_3 + k_4)$, where:
        \begin{align*}
        k_1 &= h f(t_i, w_i) \\
        k_2 &= h f(t_i + h/2, w_i + k_1/2) \\
        k_3 &= h f(t_i + h/2, w_i + k_2/2) \\
        k_4 &= h f(t_{i+1}, w_i + k_3)
        \end{align*}
        
        \textit{Step 1: Approximation at $t_1 = 0.5$}
        \begin{align*}
        k_1 &= 0.5 f(0, 0) = 0 \\
        k_2 &= 0.5 f(0.25, 0) = 0.5 \left[0.25e^{0.75}\right] \approx 0.264625 \\
        k_3 &= 0.5 f(0.25, 0.132313) = 0.5 \left[0.25e^{0.75} - 2(0.132313)\right] \approx 0.132312 \\
        k_4 &= 0.5 f(0.5, 0.132312) = 0.5 \left[0.5e^{1.5} - 2(0.132312)\right] \approx 0.988110 \\
        w_1 &= 0 + \frac{1}{6}(0 + 2(0.264625) + 2(0.132312) + 0.988110) \approx \mathbf{0.296997}
        \end{align*}
        
        \textit{Step 2: Approximation at $t_2 = 1.0$}
        \begin{align*}
        k_1 &= 0.5 f(0.5, 0.296997) \approx 0.5 [2.240845 - 0.593994] \approx 0.823425 \\
        k_2 &= 0.5 f(0.75, 0.708710) \approx 0.5 [7.115801 - 1.417420] \approx 2.849190 \\
        k_3 &= 0.5 f(0.75, 1.721592) \approx 0.5 [7.115801 - 3.443184] \approx 1.836308 \\
        k_4 &= 0.5 f(1.0, 2.133305) \approx 0.5 [20.085537 - 4.266610] \approx 7.909464 \\
        w_2 &= 0.296997 + \frac{1}{6}(0.823425 + 2(2.849190) + 2(1.836308) + 7.909464) \approx \mathbf{3.314311}
        \end{align*}
        
        \vspace{0.3cm}
        \textbf{3. Exact Solution and Comparison}
        
        We evaluate the exact solution $y(t) = \frac{1}{5}te^{3t} - \frac{1}{25}e^{3t} + \frac{1}{25}e^{-2t}$ at $t=0.5$ and $t=1.0$:
        \begin{align*}
        y(0.5) &= 0.1e^{1.5} - 0.04e^{1.5} + 0.04e^{-1} = 0.06e^{1.5} + 0.04e^{-1} \approx \mathbf{0.283617} \\
        y(1.0) &= 0.2e^{3} - 0.04e^{3} + 0.04e^{-2} = 0.16e^{3} + 0.04e^{-2} \approx \mathbf{3.219099}
        \end{align*}
        
        \textbf{Summary Comparison Table:}
        \begin{center}
        \begin{tabular}{| c | c | c | c | c | c |}
        \hline
        $t$ & Exact $y(t)$ & Mod. Euler $w_i$ & ME Error & RK4 $w_i$ & RK4 Error \\
        \hline
        0.5 & 0.283617 & 0.264625 & 0.018992 & 0.296997 & 0.013380 \\
        1.0 & 3.219099 & 3.130001 & 0.089098 & 3.314311 & 0.095212 \\
        \hline
        \end{tabular}
        \end{center}
        \textit{Note: The relatively large step size $h=0.5$ leads to significant truncation errors in both methods over the interval, particularly since the function contains rapidly growing exponential terms ($e^{3t}$).}
        \end{solution}

        \begin{problem}
            Consider the difference method
            \begin{equation*}
            w_0 = \alpha,
            \end{equation*}
            \begin{equation*}
            w_{i+1} = w_i + a_1 f(t_i, w_i) + a_2 f(t_i + \alpha_2, w_i + \delta_2 f(t_i, w_i)), \quad \text{for } i = 0, 1, \dots, N-1
            \end{equation*}
            Show that this method cannot have local truncation error $O(h^3)$ for any choice of the constants $a_1, a_2, \alpha_2$, and $\delta_2$.
            \end{problem}
            
            \solutionnewpagefalse % no newpage after the final solution
            \begin{solution}
            The local truncation error $\tau_{i+1}(h)$ is defined by replacing the approximation $w_i$ with the exact solution $y(t_i)$ and dividing the residual by $h$:
            \begin{equation*}
            \tau_{i+1}(h) = \frac{y(t_{i+1}) - y(t_i)}{h} - \frac{1}{h} \Big[ a_1 f(t_i, y(t_i)) + a_2 f(t_i + \alpha_2, y(t_i) + \delta_2 f(t_i, y(t_i))) \Big]
            \end{equation*}
            For the method to have a local truncation error of $O(h^3)$, the numerator must match the exact solution expansion up to the $O(h^3)$ terms. That is, we must have:
            \begin{equation*}
            y(t_{i+1}) - y(t_i) - \Big[ a_1 f + a_2 f(t_i + \alpha_2, y_i + \delta_2 f) \Big] = \mathcal{O}(h^4)
            \end{equation*}
            where $f$ and its derivatives are evaluated at $(t_i, y(t_i))$.
            
            \vspace{0.3cm}
            \textbf{1. Taylor Expansion of the Exact Solution}
            
            We compute the exact increment $y(t_{i+1}) - y(t_i)$ using the Taylor series of $y(t)$ around $t_i$. Since $y'(t) = f(t, y(t))$, we can find the higher derivatives using the chain rule:
            \begin{align*}
            y' &= f \\
            y'' &= f_t + f_y y' = f_t + f f_y \\
            y''' &= \frac{d}{dt}(f_t + f f_y) = (f_{tt} + f_{ty} f) + f(f_{yt} + f_{yy} f) + f_y(f_t + f_y f) \\
            &= f_{tt} + 2f f_{ty} + f^2 f_{yy} + f_t f_y + f f_y^2
            \end{align*}
            Substitute these into the Taylor series $y(t_{i+1}) = y(t_i) + h y' + \frac{h^2}{2} y'' + \frac{h^3}{6} y''' + \mathcal{O}(h^4)$:
            \begin{equation} \label{eq:exact_exp}
            y(t_{i+1}) - y(t_i) = hf + \frac{h^2}{2}(f_t + f f_y) + \frac{h^3}{6}(f_{tt} + 2f f_{ty} + f^2 f_{yy} + f_t f_y + f f_y^2) + \mathcal{O}(h^4)
            \end{equation}
            
            \vspace{0.3cm}
            \textbf{2. Taylor Expansion of the Numerical Method}
            
            Now we expand the numerical increment $T = a_1 f + a_2 f(t_i + \alpha_2, y(t_i) + \delta_2 f)$ using the multivariable Taylor series for $f$ around $(t_i, y(t_i))$:
            \begin{align*}
            f(t_i + \alpha_2, y_i + \delta_2 f) &= f + \alpha_2 f_t + \delta_2 f f_y + \frac{1}{2} \left( \alpha_2^2 f_{tt} + 2\alpha_2 \delta_2 f f_{ty} + \delta_2^2 f^2 f_{yy} \right) + \mathcal{O}(\text{order } 3)
            \end{align*}
            Multiplying by $a_2$ and adding $a_1 f$, the numerical increment is:
            \begin{equation} \label{eq:num_exp}
            T = (a_1 + a_2)f + a_2 \alpha_2 f_t + a_2 \delta_2 f f_y + \frac{1}{2} a_2 \alpha_2^2 f_{tt} + a_2 \alpha_2 \delta_2 f f_{ty} + \frac{1}{2} a_2 \delta_2^2 f^2 f_{yy} + \dots
            \end{equation}
            
            \vspace{0.3cm}
            \textbf{3. Comparison of Coefficients}
            
            For the method to be $\mathcal{O}(h^3)$, the expansions in \eqref{eq:exact_exp} and \eqref{eq:num_exp} must be identical for any arbitrary, sufficiently smooth function $f(t,y)$. 
            
            Comparing the $O(h^3)$ terms from the exact solution to the numerical method's expansion:
            \begin{itemize}
                \item The exact expansion \eqref{eq:exact_exp} contains the terms $\frac{h^3}{6} f_t f_y$ and $\frac{h^3}{6} f f_y^2$.
                \item The numerical expansion \eqref{eq:num_exp} evaluates $f$ at a single advanced point. As a result, its Taylor expansion produces purely second partial derivatives ($f_{tt}, f_{ty}, f_{yy}$) at the quadratic level. It \textbf{structurally cannot} generate products of first-order partial derivatives like $f_t f_y$ or $f f_y^2$.
            \end{itemize}
            
            Since there are no parameters $a_1, a_2, \alpha_2, \delta_2$ that can generate the required $\frac{h^3}{6} f_y(f_t + f f_y)$ terms to match the exact solution's $y'''$ component, the expansions will always differ by at least $\mathcal{O}(h^3)$. 
            
            Therefore, $y(t_{i+1}) - w_{i+1} = \mathcal{O}(h^3)$, which means dividing by $h$ leaves a local truncation error of $\tau_{i+1}(h) = \mathcal{O}(h^2)$. The method cannot have a local truncation error of $\mathcal{O}(h^3)$.
            \end{solution}

\end{document}